The \dap{} can play through Bluetooth headphones and speakers (A2DP),
which can also control playback with their own buttons (AVRCP). Only one
of them plays at a time. While one is connected, all sound goes to it.
When it goes away (it is switched off or out of range, it is disconnected,
or Bluetooth is turned off) sound goes back to the \dap{} itself:
\opt{headphone_detection}{if headphones are plugged in, playback carries
on through them; if not, the \dap{} does what
\setting{Pause on Headphone Unplug} says, as if they had just been
unplugged. With \setting{Pause and Resume}, playback then resumes when
the Bluetooth device connects again, or when headphones are plugged in.}

Devices with a display, such as car stereos, can show the title, artist
and album of the track playing, and how far into it playback is.

While a Bluetooth device plays, the volume controls set its volume, which
is separate from the headphone volume. It goes from $-127$ to 0, the
device's maximum. Devices that support it set the volume themselves, and
changing it with their own buttons shows on the \dap{} too; for others
the \dap{} lowers the sound it sends. Each device's volume is remembered
and set again when it connects. A device connecting for the first time
starts at a low volume. \setting{Volume Limit} only applies to the
headphone volume and can't be changed while a Bluetooth device plays.

\section{Bluetooth}
Turns Bluetooth \setting{On} or \setting{Off}. Bluetooth stays on, or off,
across restarts. If Bluetooth can't be started, a message says so and the
setting goes back to \setting{Off}.

\section{\label{ref:BluetoothDevices}Devices}
Lists \setting{[Pair New Device]} followed by the devices paired with
the \dap{}. The device that is connected, or being connected to, is marked
as such. If Bluetooth is off, you are asked whether to turn it on first.

\subsection{Pairing a Device}
Put the headphones or speaker in pairing mode, as its manual describes,
then select \setting{[Pair New Device]}. The \dap{} keeps searching while
the list is open and shows the audio devices it finds that are not paired
yet. Devices that don't give a name are left out.

Select a device and answer \setting{Yes} to ``Pair with \ldots?''.
\setting{No} brings you back to the list. While ``Pairing\ldots'' is shown
the \dap{} connects to the device and bonds with it; pressing
\ActionStdCancel{} gives up. Once paired, the device is connected and you
are back in \setting{Devices}. If pairing fails, a message says so and the
search goes on.

Headphones and speakers usually pair without any more questions. A device
with a display shows a number, and the \dap{} asks whether it shows the
same one: answer \setting{Yes} if it does. A device with a keypad needs the
number the \dap{} shows typed on it. Older devices are first tried with the
PIN 0000; if they refuse it, the \dap{} offers to type the device's PIN,
which its manual gives.

\subsection{Device Menu}
Selecting a paired device opens a menu titled with its name:

\begin{description}
  \item[Connect / Disconnect.] Connects to the device, or disconnects it
    if it is connected or being connected to. Nothing else is shown while
    this happens: the entry changes once the connection does, and themes
    can show the state of the connection. Connecting to a device while
    another one is connected disconnects the other one first.
  \item[Forget.] After you confirm ``Forget \ldots?'', removes the pairing
    and everything the \dap{} keeps about the device, and goes back to
    \setting{Devices}. To use the device again, pair it again. You may also
    have to remove the \dap{} from the device's own list of pairings.
  \item[Change Name.] Opens the keyboard to give the device another name
    in the \dap{}'s lists. Clearing the name brings back the device's own
    name.
  \item[Options.] Sets the options below for this device only. Until you
    set an option for a device, it follows
    \setting{Bluetooth Settings $\rightarrow$ Device Defaults}.
\end{description}

\section{Receive Audio}
Makes the \dap{} a Bluetooth speaker for a phone, tablet or computer: what
it plays comes out of the \dap{}'s own headphone or line out. If Bluetooth
is off, you are asked whether to turn it on first. Playback on the \dap{}
stops and can't be started while this screen is open, and Bluetooth
headphones are disconnected.

While the screen is open the \dap{} is visible to devices looking for one,
under the name it shows. To pair a phone with it, look for that name in the
phone's Bluetooth settings. Both show a number and the \dap{} asks whether
the phone shows the same one: answer \setting{Yes} if it does. A phone
paired once can connect again whenever the screen is open.

Once a phone is connected the screen shows its name, whether it is playing
and, if the phone tells, the title and artist. \ActionWpsPlay{} plays or
pauses on the phone, \ActionWpsSkipNext{} and \ActionWpsSkipPrev{} skip
tracks on it, and \ActionWpsVolUp{} and \ActionWpsVolDown{} set the
\dap{}'s own volume. \ActionWpsBrowse{} leaves the screen: the phone is
disconnected and the \dap{} goes back to playing
through Bluetooth headphones, connecting to the ones set to connect
automatically.

\section{\label{ref:BluetoothOptions}Device Defaults}
Each device has its own \setting{Options}, set from its menu in
\setting{Devices}. The options here are what a device uses until they are
set in its own \setting{Options}, so they also apply to newly paired
devices.

\begin{description}
  \item[Automatically Connect.] If \setting{Yes}, the \dap{} tries to
    connect to the device whenever Bluetooth is turned on, including at
    start up. Devices set this way are tried in turn until one of them
    answers. The \dap{} also lets these devices connect to it, so a pair
    of headphones switched on later can connect by itself. Choosing a
    device in the \setting{Devices} menu stops the automatic attempts.
  \item[Codec.] How sound is coded for the device. \setting{Best
    Available}, the default, gives the device the best codec it supports.
    The others set a ceiling, from plainest to best: \setting{SBC Standard
    Quality}, \setting{SBC High Quality}, \setting{aptX}, \setting{aptX
    HD} and \setting{LDAC}. All devices support SBC; the others only some.
    If the device doesn't support the codec chosen, it gets the best one
    below it that it supports, so in the end SBC. The codec is used whoever
    starts the connection: if the device sets up something else, it is asked
    to switch.

    \setting{SBC High Quality} uses up to about 330~kbit/s and
    \setting{SBC Standard Quality} up to about 230~kbit/s, which is easier
    to keep up with when the radio connection is poor. Either way, Rockbox
    lowers the SBC rate on its own while the connection can't keep up.
    The other codecs need more processing power. A change takes effect the
    next time the device connects.
\end{description}

\section{Stream at 44.1 kHz}
If \setting{Yes}, everything sent to a device is resampled to 44.1~kHz, the
rate of most music. If \setting{No}, a track at another rate the device
supports, such as 48~kHz, is sent at that rate. The same bitrate at 48~kHz
takes more packets a second, which a poor radio connection may not carry,
so \setting{Yes} is the default. A change takes effect the next time a
device connects.

\note{Whether Bluetooth is on, \setting{Stream at 44.1 kHz} and the
  \setting{Device Defaults} are kept
  with the other settings in \fname{/.rockbox/config.cfg}. The pairings and
  everything set for each device are kept in
  \fname{/.rockbox/bt\_config.conf}, one section per device. Device
  settings are stored as \config{RbAlias}, \config{RbAutoConnect} and
  \config{RbCodec} entries next to the pairing keys. If the Bluetooth chip
  has no address of its own, the \dap{} makes one up when Bluetooth is
  first turned on and keeps it in \fname{/.rockbox/bt\_addr.txt}.}
